\section{Описание}
Требуется реализовать программную библиотеку для словаря на основе **AVL-дерева** — самобалансирующегося бинарного дерева поиска. AVL-дерево гарантирует высоту $\log n$ и время выполнения операций вставки, удаления и поиска $O(\log n)$ за счёт автоматической балансировки после каждой операции (разница высот левого и правого поддеревьев не превышает 1).

Для ключей (регистронезависимые строки английских букв длиной $\le 256$) сравнение и хранение выполняются в нижнем регистре. Значением является 64-битное беззнаковое целое число.

Реализация содержит шаблонный класс \texttt{avl::AVL} с необходимыми поворотами (left, right, left-right, right-left) и вспомогательные функции для удаления.

\begin{lstlisting}[language=C++]
class AVL {
   private:
    struct Node {
        Key key;
        Value value;
        int height = 1;
        Node* left = nullptr;
        Node* right = nullptr;

        Node(const Key& k, const Value& v) : key(k), value(v) {}
    };

    Node* root_ = nullptr;
    size_t size_ = 0;

    int get_height(Node* node) const { return node ? node->height : 0; }

    int balance_factor(Node* node) const {
        return node ? get_height(node->right) - get_height(node->left) : 0;
    }

    void fix_height(Node* node) {
        if (!node) return;
        int hl = get_height(node->left);
        int hr = get_height(node->right);
        node->height = (hl > hr ? hl : hr) + 1;
    }

    Node* rotate_right(Node* p) {
        Node* q = p->left;
        p->left = q->right;
        q->right = p;
        fix_height(p);
        fix_height(q);
        return q;
    }

    Node* rotate_left(Node* q) {
        Node* p = q->right;
        q->right = p->left;
        p->left = q;
        fix_height(q);
        fix_height(p);
        return p;
    }

    Node* balance(Node* p) {
        if (!p) return nullptr;
        fix_height(p);

        if (balance_factor(p) == 2) {
            if (balance_factor(p->right) < 0) p->right = rotate_right(p->right);
            return rotate_left(p);
        }
        if (balance_factor(p) == -2) {
            if (balance_factor(p->left) > 0) p->left = rotate_left(p->left);
            return rotate_right(p);
        }
        return p;
    }

    Node* insert(Node* node, const Key& key, const Value& value,
                 bool& success) {
        if (!node) {
            success = true;
            ++size_;
            return new Node(key, value);
        }

        if (key < node->key)
            node->left = insert(node->left, key, value, success);
        else if (key > node->key)
            node->right = insert(node->right, key, value, success);
        else
            success = false;

        return balance(node);
    }

    Node* remove_min(Node* node, Node*& min_node) {
        if (!node->left) {
            min_node = node;
            return node->right;
        }
        node->left = remove_min(node->left, min_node);
        return balance(node);
    }

    Node* remove(Node* node, const Key& key, bool& success) {
        if (!node) {
            success = false;
            return nullptr;
        }

        if (key < node->key)
            node->left = remove(node->left, key, success);
        else if (key > node->key)
            node->right = remove(node->right, key, success);
        else {
            success = true;
            --size_;

            if (!node->left || !node->right) {
                Node* temp = node->left ? node->left : node->right;
                delete node;
                return temp;

            } else {
                Node* min_node = nullptr;
                node->right = remove_min(node->right, min_node);

                min_node->left = node->left;
                min_node->right = node->right;

                delete node;
                return balance(min_node);
            }
        }
        return balance(node);
    }

    // ... etc.

    void serialize(Node* node, std::ofstream& out) const {
        if (!node) return;

        uint16_t len = static_cast<uint16_t>(node->key.size());
        out.write(reinterpret_cast<const char*>(&len), sizeof(len));
        out.write(node->key.data(), len);
        out.write(reinterpret_cast<const char*>(&node->value),
                  sizeof(node->value));

        serialize(node->left, out);
        serialize(node->right, out);
    }

    // ... etc.
}
\end{lstlisting}
\pagebreak